\section*{SN-RRF Continuum-Family Synthesis Addendum}

\subsection*{Summary}

This addendum extends the earlier A10/TVC-RRF and SN-RRF diagnostic sequence by replacing fixed-exponent language with a continuum-family interpretation. The original A10/TVC-derived kernel used a locked exponent $n = 3.258993$. Subsequent curve-level null-refit controls showed that this locked kernel defeats destructive nulls, but neighboring monotone-saturating kernels compete. Stage 30 therefore scanned a continuum grid of saturating kernels,
\[
K_n(x) = \frac{x^n}{1+x^n},
\]
across 41 values from $n=0.5$ to $n=12.0$.

The Stage 30 continuum audit covered 143 valid galaxies with no missing or unreadable Rotmod cases. The best-$n$ median was 2.25, with an interquartile range of [1.5, 3.2544965]. The locked exponent was within $+2$ AIC of the best continuum exponent for 0.322 of galaxies and within $+10$ AIC for 0.636 of galaxies. The median $\Delta$AIC(locked $n$ minus best continuum $n$) was 6.292. All six Stage 30 guards passed.

\subsection*{Interpretation}

The continuum profile does not support the uniqueness of $n = 3.258993$ as a universal fixed exponent. Instead, it supports a heterogeneous monotone-saturating residual-response family. In this corrected language, the A10/TVC exponent is retained as a representative and historically motivated diagnostic anchor, not as a physical law.

\subsection*{Claim Boundary}

This addendum does not claim dark-matter exclusion, Lambda-CDM replacement, MOND/RAR defeat, a Bullet-Cluster explanation, proof of a TVC mechanism, or a Hubble-tension solution. It supports only a diagnostic continuum-family interpretation of SPARC rotation-curve residual responses under the tested protocol.
